\documentclass[11pt,a4paper]{article}

\usepackage[utf8]{inputenc}
\usepackage[T1]{fontenc}
\usepackage{amsmath, amsfonts, amssymb}
\usepackage{geometry}
\usepackage{graphicx}
\usepackage{longtable}
\usepackage{pdfpages}
\usepackage{url}
\usepackage{microtype}

\geometry{margin=1in}

\title{\textbf{Nichols Radial Injection Model (RIM) IV:\\
A 5D Mass-Regulated Manifold Solution to Cosmic Expansion}}

\author{Lawrence William Nichols}
\date{28 February 2026}

\begin{document}
\maketitle

\begin{abstract}
This paper presents the Nichols Radial Injection Model (RIM), a cosmological framework that replaces the Dark Energy placeholder with the mechanical displacement of a 5D hypersphere manifold. By defining a fundamental curvature constant $\Omega_k = 0.004$, we identify a True Hypersphere Radius ($R$) of 735 Gly. We demonstrate a ``System Architecture'' lock where independent metrics—including the 46.5 Gly CMB skin ratio and the 1.19 km/s resonant pulse speed—converge on this radius. The model suggests the universe is a mass-regulated pressure vessel governed by a recursive scaling constant of 0.1333.
\end{abstract}

\section{Introduction}
The RIM performs a mechanical audit of cosmic expansion. Unlike conventional cosmology, which relies on Dark Energy, the universe is treated as a 5D mass-regulated hypersphere with measurable pulse dynamics. The framework allows direct calculation of hypersphere radius, pulse speed, and expansion harmonics.

\section{Fundamental Curvature and Hypersphere}
The RIM defines the \textbf{Curvature Constant} as the physical track of the manifold:
\begin{equation}
k = 0.004
\end{equation}

Given the observable radius ($r$) of 46.5 billion light-years, the True Hypersphere Radius ($R$) is:
\begin{equation}
R = \frac{46.5}{\sqrt{0.004}} \approx 735 \ \text{Gly}
\end{equation}

\section{Recursive Scaling and Stability}
The architecture follows a recursive scaling constant of \textbf{0.1333}. The relationship between the macro-scale radius and the fundamental curvature is locked by the system architecture:
\begin{equation}
735 \times (0.1333)^6 \approx 0.004
\end{equation}
This confirms that the curvature observed at the light horizon is the same "code" that governs the smallest injection point.

\section{Historical Pulse Timeline (1607--2025)}
\begin{longtable}{|c|p{5cm}|p{5cm}|}
\hline
\textbf{Year (Pulse)} & \textbf{Event / Pulse / Solar / Cosmic} & \textbf{Volcano / Notes} \\ \hline
1607 & First Telescope Observations (Galileo / Kepler) & 1606 --- No major VEI-5+ (control) \\ \hline
1618 & Triple Comet / Kepler's 3rd Law & 1617 --- Huaynaputina (1600) VEI-6 \\ \hline
1629 & Pulse & --- \\ \hline
1640 & Pulse & 1639 --- Komaga-take (1640) major eruption \\ \hline
1651 & Pulse & --- \\ \hline
1662 & Pulse & 1661 --- Vesuvius (1660) direct hit \\ \hline
1673 & Pulse & --- \\ \hline
1684 & Pulse & --- \\ \hline
1695 & Pulse & 1694 --- Santorini (1707--1711) eruption sequence \\ \hline
1706 & Maunder Minimum era & --- \\ \hline
1717 & Pulse & 1716 --- Taal (1716) direct hit \\ \hline
1728 & Pulse & 1727 --- Lanzarote (1730--1736) \\ \hline
1739 & Pulse & --- \\ \hline
1750 & Pulse & 1749 --- Katla (1755) major Iceland eruption \\ \hline
1761 & Transit of Venus & 1760 --- Laki (1783) VEI-6 in window \\ \hline
1772 & Cycle 2 Peak & --- \\ \hline
1783 & Cycle 3 Peak & 1760 --- Laki (1783) close alignment \\ \hline
1794 & Pulse & 1793 --- Unzen (1792) deadliest Japan eruption \\ \hline
1805 & Battle of Trafalgar / Cycle 5 Peak & 1804 --- St.\ Helens (1800), Taal (1808) \\ \hline
1816 & Dalton Minimum Peak & 1815 --- Tambora VEI-7, direct hit \\ \hline
1838 & Stellar Parallax (61 Cygni) & 1848 --- Cosig\"uina (1835) VEI-5/6 \\ \hline
1849 & First Photo of the Sun & Continued window \\ \hline
1860 & Carrington Event / Eta Carinae eruption & 1859 --- Mauna Loa (1859) \\ \hline
1871 & Great Chicago Fire & 1870 --- Mauna Loa cluster \\ \hline
1882 & Great September Comet & 1881 --- Tarawera (1886) \\ \hline
1893 & Great Sunspot & 1892 --- Krakatoa (1883) VEI-6 \\ \hline
1904 & Mount Wilson Observatory & 1903 --- Santa María (1902) VEI-6 \\ \hline
1915 & Einstein publishes GR & 1914 --- Sakurajima VEI-4/5 \\ \hline
1926 & Hubble spike ($H_0\approx78$) & 1925 --- Rabaul activity \\ \hline
1937 & Cosmic Ray Neutrons & 1936 --- Aleutian activity \\ \hline
1948 & Palomar 200-inch Telescope & 1947 --- Hekla eruption \\ \hline
1959 & Modern Maximum / Cycle 19 & 1958 --- Kīlauea activity spike \\ \hline
1970 & Apollo Era / Cycle 20 peak & 1969 --- Deception Island eruptions \\ \hline
1981 & First Space Shuttle Flight / Cycle 21 & 1980 --- Mount St.\ Helens VEI-5 \\ \hline
1992 & COBE mission / Cycle 22 peak & --- \\ \hline
2003 & Halloween Storms / Cycle 23 & 1991 --- Pinatubo VEI-6 \\ \hline
2014 & Double Peak Solar Cycle 24 & --- \\ \hline
2025 & GRB 250702B / Cycle 25 peak & --- \\ \hline
\end{longtable}

\section{Hypersphere Parameters and Constants}
\begin{longtable}{|p{3.5cm}|c|c|p{3.5cm}|}
\hline
\textbf{Parameter} & \textbf{Value} & \textbf{Unit} & \textbf{Significance} \\ \hline
True Hypersphere Radius ($R$) & 735 & Gly & Fundamental 5D spatial radius \\ \hline
Curvature Constant ($k$) & 0.004 & --- & The Physical "Track" \\ \hline
Pulse Speed ($V_p$) & 1.19 & km/s & Constant Mechanical Drive \\ \hline
Observable Skin ($r$) & 46.5 & Gly & Light horizon; Harmonic ratio $\approx 0.063$ \\ \hline
Scaling Constant & 0.1333 & --- & Recursive system architecture \\ \hline
Corridor Width & 10 & Gly & Expansion Rail / Safety Buffer \\ \hline
Pulse Period & 11.07 & Years & Resonant frequency \\ \hline
Total Circumference & 4.6 & tLy & Hypersphere outer boundary \\ \hline
Universal Mass & $1.5\times10^{53}$ & kg & Total cosmic mass \\ \hline
Surface Age & 16.6 & Gyr & Current injection duration \\ \hline
Bulk Age & $\sim4.5\times10^{12}$ & Years & Substrate age \\ \hline
\end{longtable}

\section{Conclusion}
The Nichols Radial Injection Model replaces Dark Energy with a mechanical 5D hypersphere framework. The 0.004 curvature is the proof of the track.

\section{Richard Feynman}
It doesn't matter how beautiful your theory is, it doesn't matter how smart you are. If it doesn't agree with experiment, it's wrong.

\begin{thebibliography}{9}
\bibitem{gompertz2025}
Gompertz, B.\ P., et al.,
``JWST Spectroscopy of GRB 250702B,''
arXiv:2509.22778, 2025.
\url{https://arxiv.org/abs/2509.22778}

\bibitem{sn1054_wiki}
Wikipedia,
``SN 1054,''
2026.
\url{https://en.wikipedia.org/wiki/SN_1054}

\bibitem{kepler1604_nasa}
NASA,
``420 Years Ago: Astronomer Johannes Kepler Observes a Supernova,''
2024.
\url{https://www.nasa.gov/history/420-years-ago-astronomer-johannes-kepler-observes-a-supernova}

\bibitem{einstein1935}
Einstein, A., \& Rosen, N. (1935).
\textit{The Particle Problem in the General Theory of Relativity}.
Physical Review, 48(1), 73.
\end{thebibliography}
\newpage

\section{Engineering Studies and Visual Data}
This section includes relevant visual diagrams and engineering models.
\includepdf[pages=-]{(RIM) IV.pdf}

\end{document}
